\documentclass[25pt,
a0paper,
portrait,
svgnames,
blockverticalspace=0.3cm
]{tikzposter}
%\usepackage{extsizes}
\usepackage{xcolor}
\usepackage{graphicx}
\usepackage{listings}
\usepackage{listingsutf8}
\usepackage{verbatim}
\usepackage{fancyvrb}
\usepackage{amsmath,amsfonts,amsthm,bm} % Math packages
\usepackage{CedilleVerbatim}
\usepackage{url}


% lstlisting coq style (inspired from a file of Assia Mahboubi)
%
\definecolor{myred}{HTML}{c82829} 
\definecolor{myorange}{HTML}{f5871f} 
\definecolor{myyellow}{HTML}{eab700} 
\definecolor{mygreen}{HTML}{718c00} 
\definecolor{myaqua}{HTML}{3e999f} 
\definecolor{myblue}{HTML}{4271ae} 
\definecolor{mypurple}{HTML}{8959a8} 
\lstdefinelanguage{Coq}{ 
%
% Anything betweeen $ becomes LaTeX math mode
mathescape=true,
moredelim=[is][\color{red}]{@}{@},
%
% Comments may or not include Latex commands
texcl=false, 
%
% Vernacular commands
morekeywords=[1]{Section, Module, End, Require, Import, Export,
  Variable, Variables, Parameter, Parameters, Axiom, Hypothesis,
  Hypotheses, Notation, Local, Tactic, Reserved, Scope, Open, Close,
  Bind, Delimit, Definition, Let, Ltac, Fixpoint, CoFixpoint, Add,
  Morphism, Relation, Implicit, Arguments, Unset, Contextual,
  Strict, Prenex, Implicits, Inductive, CoInductive, Record,
  Structure, Canonical, Coercion, Context, Class, Global, Instance,
  Program, Infix, Theorem, Lemma, Corollary, Proposition, Fact,
  Remark, Example, Proof, Goal, Save, Qed, Defined, Hint, Resolve,
  Rewrite, View, Search, Show, Print, Printing, All, Eval, Check,
  Projections, inside, outside, Def},
%
% Gallina
morekeywords=[2]{forall, exists, exists2, fun, fix, cofix, struct,
  match, with, end, as, in, return, let, if, is, then, else, for, of,
  nosimpl, when},
%
% Sorts
morekeywords=[3]{Type, Prop, Set, option},
%
% Various tactics, some are std Coq subsumed by ssr, for the manual purpose
morekeywords=[4]{pose, set, move, case, elim, apply, clear, hnf,
  intro, intros, generalize, rename, pattern, after, destruct,
  induction, using, refine, inversion, injection, rewrite, congr,
  unlock, compute, ring, field, fourier, replace, fold, unfold,
  change, cutrewrite, simpl, have, suff, wlog, suffices, without,
  loss, nat_norm, assert, cut, trivial, revert, bool_congr, nat_congr,
  symmetry, transitivity, auto, split, left, right, autorewrite},
%
% Terminators
morekeywords=[5]{by, done, exact, reflexivity, tauto, romega, omega,
  assumption, solve, contradiction, discriminate},
%
% Control
morekeywords=[6]{do, last, first, try, idtac, repeat},
%
% Comments delimiters, we do turn this off for the manual
morecomment=[s]{(*}{*)},
%
% Spaces are not displayed as a special character
showstringspaces=false,
%
% String delimiters
morestring=[b]",
morestring=[d]’,
%
% Size of tabulations
tabsize=3,
%
% Enables ASCII chars 128 to 255
extendedchars=false,
%
% Case sensitivity
sensitive=true,
%
% Automatic breaking of long lines
breaklines=false,
%
% Default style fors listings
basicstyle=\fontfamily{monospace}\selectfont\small,
%
% Position of captions is bottom
captionpos=b,
%
% flexible columns
columns=[l]flexible,
%
% Style for (listings') identifiers
identifierstyle={\ttfamily\color{black}},
% Style for declaration keywords
keywordstyle=[1]{\ttfamily\color{myblue}},
% Style for gallina keywords
keywordstyle=[2]{\ttfamily\color{mypurple}},
% Style for sorts keywords
keywordstyle=[3]{\ttfamily\color{mygreen}},
% Style for tactics keywords
keywordstyle=[4]{\ttfamily\color{myaqua}},
% Style for terminators keywords
keywordstyle=[5]{\ttfamily\color{DarkRed}},
%Style for iterators
%keywordstyle=[6]{\ttfamily\color{Darkpink}},
% Style for strings
stringstyle=\ttfamily,
% Style for comments
commentstyle={\ttfamily\color{DarkGreen}},
%
%moredelim=**[is][\ttfamily\color{red}]{/&}{&/},
literate=
    {\\forall}{{\color{DarkGreen}{$\forall\;$}}}1
    {\\exists}{{$\exists\;$}}1
    {<-}{{$\leftarrow\;$}}1
    {=>}{{$\Rightarrow\;$}}1
    {==}{{\code{==}\;}}1
    {==>}{{\code{==>}\;}}1
%    {:>}{{\code{:>}\;}}1
    {->}{{$\rightarrow\;$}}1
    {<->}{{$\leftrightarrow\;$}}1
    {<==}{{$\leq\;$}}1
    {\#}{{$^\star$}}1 
    {\\o}{{$\circ\;$}}1 
    %{\@}{{$\cdot$}}1 
    {\/\\}{{$\wedge\;$}}1
    {\\\/}{{$\vee\;$}}1
    {++}{{\code{++}}}1
    {~}{{$\sim$}}1
    {\@\@}{{$@$}}1
    {\\mapsto}{{$\mapsto\;$}}1
    {\\hline}{{\rule{\linewidth}{0.5pt}}}1
%
}[keywords,comments,strings]

\lstnewenvironment{coq}{\lstset{language=Coq}}{}

% pour inliner dans le texte
\def\coqe{\lstinline[language=Coq, basicstyle=\normalsize]}
% pour inliner dans les tableaux / displaymath...
\def\coqes{\lstinline[language=Coq, basicstyle=\scriptsize]}

%%% Local Variables: 
%%% mode: latex
%%% Local IspellDict: british
%%% TeX-master: "main.tex"
%%% End: 


\newcommand{\inlinecoq}[1]{\lstinline[breaklines=true, language=Coq]|#1|}
\tikzposterlatexaffectionproofoff

\let\thempfootnote\thefootnote%http://tex.stackexchange.com/questions/956/footnotemark-and-footnotetext-in-minipage#959
\newcommand\printfootnote[1]{
\addtocounter{footnote}{1}%
\footnotetext{#1}}

\title{How small can we make a useful type theory?}
\author{Pedro da Costa Abreu Jr.}
\institute{Purdue University, USA \\ pdacost@purdue.edu}

% \usetheme{Desert}
\usetheme{Simple}

\begin{document}

\maketitle

% \node[anchor=east] at (TP@title.east) {\includegraphics[width=0.20\linewidth]{images/purdue-logo.png}};

% \block{The One Million Dollar Question}{
% \fontsize{25pt}{60pt}\selectfont

\begin{columns}

\column{.45}
\block{There are too many posters here, why should I read yours?}{
\textbf{Universe Hierarchy}, \textbf{datatypes} and \textbf{large elimination} are well known features 
of modern theorem provers. However, they are also known to
\textbf{complicate metatheorical properties}, such as soundness, type inference, and
normalization.

This work aims to assess how much these features
are actually necessary in practice. To achieve this we build
\textbf{Coquedille, a compiler from a subset of Gallina to Cedille\textsuperscript{1}}, and
probe it by translating as much as possible of the Software Foundations Series.

% The main goal of this work is to determine how much these features are actually useful
% in practice. We aim to achieve this by building \textbf{Coquedille, a compiler
% from Coq to Cedille}.
}
\block{What do you mean Universe Hierarchy?}{
Recall Russell's paradox: 
Imagine the set of all sets that does not include itself
$$ \Delta = \{ x \ | \ x \notin x \}$$
The million dollar question is: $\Delta \in \Delta$?

If $\Delta \in \Delta$ then $\Delta \notin \Delta$,
conversely, if $\Delta \notin \Delta$ then $\Delta \in \Delta$!.
This contradiction indicates that such $\bm{\Delta}$ \textbf{simply cannot exist}.

This means that a sound theory needs to disallow such definitions. \textbf{One fix to this problem is to only allow definitions to quantify over terms smaller than what is being defined.}
This gives rise to a \textit{hierarchy} of larger and larger types.
In Coq this is Prop, Set, Type(1), Type(2), ...

Another fix to this problem is to disallow these definitions altogether, resulting in a weaker theory.
}

\block{Cool, what about Large Eliminations?}{
\textbf{Large Eliminations are function that computes types from terms.}
The canonical example is computing a proposition from an arbitrary term, such as:
\lstinputlisting[language=Coq]{bool_to_prop.v}
See right for an example of these in action.
}

\block{By Datatypes you mean Algebraic Datatypes, right?}{
Yep! Here are two cherrypicked examples for you:

\begin{tabular}{ll}
    \begin{minipage}[t]{.5\linewidth}
        \lstinputlisting[language=Coq]{vec.v}
    \end{minipage}
    &
    \begin{minipage}[t]{.5\linewidth}
        \lstinputlisting[language=Coq]{eq.v}
    \end{minipage}
\end{tabular}

}

 \block{So you are saying Cedille has none of this?}{
 Exactly! Cedille has \textbf{no large eliminations}, and \textbf{no primitive datatypes}, and like Haskell and Agda,
 \textbf{no universe universe hierarchy}, all of this implemented in an extrinsic type theory.\\
 
 \textit{Ask me how can we get away without datatypes.}
 }
 
 
\block{Let's talk about Coquedille}{
  Of course, that's why we're here for!

  Coquedille is a translater from Coq to Cedille.
  \textbf{We implement Coquedille in Coq} itself, this allows us to \textbf{use metacoq} to access the abstract syntax tree of the gallina terms.

  Since Coq and Cedille are based on the Calculus of Constructors, most of the translation is straightforward!


    % here we talk about useful. How far we can get with these features. What is our vehicle?

    % Well, people like Coq, so we can probe existing coq developing if they have these stuff, we have started implementing coquedille to answer this question.
}


% \block{Wait a second... Extrinsic Type Theory?}{
% Oh right, I knew I was forgeting something.
% There is one more important mismatch between Coq and Cedille: \textbf{Coq is an Intrinsic Type theory},
% but \textbf{Cedille is an Extrinsic Type Theory}.

% The key difference is that in an Intrinsic Type Theory we would write $\bm{\lambda x: T . e}$,
%  while in an Extrinsic Type Theory, it would be $\bm{\lambda x . e}$.
% I know, it seems like a minute difference but it leads to huge consequences, e.g. type uniqueness and decidability of type checking.
% In any case, this mismatch opens up the interesting question of how these two flavours of type theories will
% interact with one another soundly.
% }


\column{.03}

\column{.47}
\block{Ok... What can you do with it?}{


\begin{tabular}[t]{p{.6\linewidth}r}
    \begin{minipage}[t]{.9\linewidth}

      Here is an example for you: let's \textbf{translate the following non-trivial proof:
      List nil and Vector vnil are not the same}.

    Since we are comparing terms of two different types we cannot use
    the regular equality, so we use the JMeq heterogeneous equality instead (notated as {\small'$\sim=$'}).

    Don't waste your time trying to understand this proof! All I want you to notice is that 
    \textbf{the highlighted part is a large elimination}.\\

    \textit{If you want to see the proof written with tactics, take a look at the project github\textsuperscript{2}}

    \end{minipage}
    &
    \begin{minipage}[t]{0.50\linewidth}
    \lstinputlisting[language=Coq]{nil_neq_vnil.v}
    \end{minipage}
\end{tabular}

}


 \block{Hey! You just said Cedille doesn't support Large Eliminations!}{
 Preciselly! How can we translate it then?
 We use a nice trick of translating it to
 a Cedille term $\bm{\delta}$, which basically 
 \textbf{checks if two lambda terms are \bm{$\beta\eta$}-equivalent} via the
 Bohm-out algorithm\textsuperscript{3}. If the terms being compared are not $\beta\eta$-equivalent, $\delta$
 allows us to apply the explosion principle
 %redundancy here is intentional%
 to prove anything.
 %------------------------------%
 In order to do this translation, we first prove a lemma
 relating Coq equality \lstinline[language=Coq, basicstyle=\large]{eq} to Cedille
 built-in equality \Verb[fontsize=\normalsize]+\{x ≃ y\}+\\
 
 \VerbatimInput[fontsize=\small,commandchars=+\^\&]{eqprimeq.ced}
 
Now we can translate the large elimination by
using $\delta$ on the generated proof that \lstinline[language=Coq, basicstyle=\large]{nil = cons a nil} (look at the type of H4 at the proof above).\\
\VerbatimInput[fontsize=\small,commandchars=+\^\&]{cedille-proof.ced}
}

\block{Next Steps}{
\begin{itemize}
    \item Translate induction and recursion
    \item Check how much of Software Foundations we are able
    to translate even with these shortcomings
    \item Prove the correctness of the translation
    \begin{itemize}
      \item Translation Validation
      \item Logical Relations
      \item ???
    \end{itemize}
\end{itemize}
}

\block{Wait, I still have a question}{
  I'm more than happy to hear any insights, doubts, feedbacks, or I don't know, ``hi''s.
}

% \block{Shameless Plug}{
% I'm looking for summer internships!
% }

\block{Citations}{
\textsuperscript{1} {\mbox{\url{https://github.com/cedille/cedille}}}.

\textsuperscript{2} {\mbox{\url{https://github.com/pedrotst/coquedille}}}.

\textsuperscript{3} H. Barendregt, The Lambda Calculus; Its Syntax and Semantics.


}

\end{columns}

\end{document}

%%% Local Variables:
%%% mode: latex
%%% TeX-master: t
%%% End:
